% Half title, title page and colophon.
%
% \bookversion is the release the PDF was built for: `make VERSION=v2` passes it in, and the
% release workflow passes in the tag. Without it the title page carries only the build date.
\providecommand{\bookversion}{}
\newcommand{\bookdate}{\number\day\space\ifcase\month\or January\or February\or March\or April\or
  May\or June\or July\or August\or September\or October\or November\or December\fi\space
  \number\year}
\newcommand{\bookedition}{\ifx\bookversion\empty\else\bookversion\,·\,\fi\bookdate}

\thispagestyle{empty}
\vspace*{2.2in}
\begin{flushright}
  {\Large\bfseries Communication Protocols and Drivers}
\end{flushright}
\cleardoublepage

\thispagestyle{empty}
\vspace*{1.1in}
\begin{flushleft}
  {\color{accent}\rule{0.9in}{2.5pt}}\par\vspace{1.4ex}
  {\fontsize{30}{36}\selectfont\bfseries Communication\\[0.15ex] Protocols and\\[0.15ex]
   Drivers\par}
  \vspace{2.2ex}
  {\Large\itshape From a protocol of your own to a\\tested CAN driver in C++\par}
  \vspace{1.0in}
  {\large Erik Pihl\par}
  \vfill
  {\small\color{muted} \bookedition\par}
\end{flushleft}
\newpage

\thispagestyle{empty}
\vspace*{\fill}
{\small
\noindent\textit{Communication Protocols and Drivers}\\
Course book. One chapter per lecture.

\medskip
\noindent The book takes you from a protocol of your own, where every mechanism is written by
hand, to a driver for a CAN controller in an FPGA that a parallel class is designing at the same
time. The last third connects the two halves on real hardware.

\medskip
\noindent This is the English edition of \textit{Kommunikationsprotokoll och drivrutiner}. Both
editions are set from the same course material, and the two are kept in step chapter by chapter.

\medskip
\noindent Text and figures: Erik Pihl.

\medskip
\noindent The text and figures of this book, like the course material they are set from, are
licensed under the Creative Commons Attribution 4.0 International licence (CC BY 4.0). You may
copy, distribute and adapt the material for any purpose, including commercially, provided that
you credit the author and the licence. The licence is at
\url{https://creativecommons.org/licenses/by/4.0/}.
\medskip
\noindent The code in the book's listings may also be used under the MIT licence, which is the
licence for all code in the course repository. The lecture material and the suggested solutions
to the exercises live there too.

\medskip
\noindent The code is C++17. The driver keeps to a subset that can be cross-compiled freestanding
for an AVR32DB28, that is, without heap, exceptions and RTTI. AVR, AVR-Dx, Cyclone V and DE0-CV
are trademarks of their respective owners.

\medskip
\noindent Set with Lua\LaTeX\ in TeX Gyre Pagella and DejaVu Sans Mono.

\medskip
\noindent\textcolor{muted}{This edition: \bookedition.}
\par}
\cleardoublepage
